\documentclass[11pt,a4paper]{article}
\usepackage[T1]{fontenc}
\usepackage[utf8]{inputenc}
\usepackage{lmodern}
\usepackage[margin=1in]{geometry}
\usepackage{amsmath,amssymb,amsthm,booktabs}
\usepackage{microtype}
\usepackage[numbers,sort&compress]{natbib}
\usepackage[hidelinks]{hyperref}\hypersetup{pdftitle={Domination log-concavity for the caterpillars T(p,0,0,q)},pdfauthor={Xiaojian Zeng},pdfsubject={Log-concavity for a two-parameter caterpillar family}}
\newtheorem{theorem}{Theorem}
\newtheorem{lemma}[theorem]{Lemma}
\newcommand{\ind}{\mathbf{1}}
\newcommand{\LC}{\operatorname{LC}}
\title{Domination log-concavity for the caterpillars\\$T(p,0,0,q)$}
\author{Xiaojian Zeng\\[3pt]
  \small PBC School of Finance, Tsinghua University, 100190\\
  \small\texttt{xiaojianzeng@foxmail.com}}
\date{}

\begin{document}
\maketitle

\begin{abstract}
We prove that the domination polynomial of the caterpillar $T(p,0,0,q)$
has a log-concave coefficient sequence for all integers $p,q\geq1$.
This family has four spine vertices, with leaves attached only at the two
endpoints, and lies outside the known all-positive-leg theorem.
An endpoint-state decomposition, followed by reversal, expresses the polynomial
through binomial convolutions and three sparse additions. The additions affect
only ten log-concavity differences, for which we give explicit nonnegative
coefficient certificates. The general claim for caterpillars with arbitrary
zero-leg patterns remains unresolved by this result.
\end{abstract}

\section{The restricted problem}

For a finite simple graph $G$, a set $W\subseteq V(G)$ is \emph{dominating}
if every vertex belongs to $W$ or has a neighbor in $W$. Write
\[
  D_G(x)=\sum_{j=0}^{|V(G)|}d_j(G)x^j,
  \qquad d_j(G)=\#\{W: W\text{ is dominating and }|W|=j\}.
\]
This is the domination polynomial; see Kotek, Preen, Simon, Tittmann,
and Trinks~\cite[Definition~1.1]{Kotek2012}.
A nonnegative sequence $(c_0,\ldots,c_n)$ is \emph{log-concave} if
$c_j^2\geq c_{j-1}c_{j+1}$ for $1\leq j\leq n-1$.
All coefficients, including initial zeros, are included in this convention.

Let $T(a_1,\ldots,a_k)$ consist of a path $v_1\cdots v_k$, called its
spine, with $a_i$ additional leaves attached to $v_i$.
Du, Heilman, and Panova~\cite[Theorem~3.3]{Du2026} prove log-concavity when
every $a_i\geq1$, using the factorization
\begin{equation}\label{eq:positive}
  D_{T(a_1,\ldots,a_k)}(x)
    =\prod_{i=1}^k\bigl(x^{a_i}+x(1+x)^{a_i}\bigr).
\end{equation}
Their Remark~3.4 expresses the belief that the conclusion also holds when
some $a_i=0$. The corresponding universal question, for every $k\geq1$
and every nonnegative integer vector $(a_1,\ldots,a_k)$, remains unresolved
here. Our contribution is the following proper subfamily.

\begin{theorem}\label{thm:main}
For all integers $p,q\geq1$, let $T=T(p,0,0,q)$ and $N=p+q+4$.
Then
\[
  d_j(T)^2\geq d_{j-1}(T)d_{j+1}(T)
  \qquad(1\leq j\leq N-1).
\]
\end{theorem}

The role of positive legs can also be described through leaf components.
In the terminology used by Zeng~\cite[Section~2]{Zeng2026}, the internal
core is induced by the nonleaf vertices, and a tree is proper if each core
vertex is adjacent to a leaf. Deleting the internal edges of a proper tree
leaves stars, each of order at least two. For $T(p,0,0,q)$ the same deletion
instead gives component orders $p+1,1,1,q+1$: the two singleton components
are precisely the bare spine vertices. This shared core description explains
the structural change from the all-positive case; the cited reconstruction
results concern the chromatic symmetric function, a different invariant.

In~\eqref{eq:positive}, an unselected spine vertex is dominated by its
forced leaves. A bare vertex has no such guarantee, so its state couples to
its spine neighbors. We resolve that coupling explicitly for two consecutive
bare vertices. We then use the standard convolution principle
\cite[Section~2]{Gross2015} and check the few differences altered by the
remaining additions. The cited caterpillar theorem does not contain this case.
Our bounded literature search through 12 September 2026 found no treatment of
this precise family; this is not a claim of exhaustive priority.

\section{Polynomial decomposition and log-concavity}

Put $r=p+q$. If the left endpoint of the spine is selected, it and its
leaves contribute $U_p=x(1+x)^p$. If it is unselected, every leaf must be
selected, giving $V_p=x^p$; these leaves also dominate the endpoint because
$p\geq1$. Define $U_q,V_q$ at the right endpoint in the same way.
When both endpoints are selected, the two middle vertices can be chosen
freely. In each other endpoint state, at least one middle vertex must be
selected: without one, a middle vertex next to an unselected endpoint is
undominated. Conversely, selecting either middle vertex dominates both of
them. Thus their contribution is $(1+x)^2$ in the first state and
$(1+x)^2-1=x(2+x)$ in the other three states. This gives
\[
 D_T(x)=U_pU_q(1+x)^2+(U_pV_q+V_pU_q+V_pV_q)x(2+x).
\]
Consequently $D_T(x)=x^2H(x)$, where
\begin{equation}\label{eq:H}
\begin{split}
 H(x)={}&(1+x)^{r+2}
       +x^p(2+x)(1+x)^q+x^q(2+x)(1+x)^p\\
       &+x^{r-1}(2+x).
\end{split}
\end{equation}
The degree of $H$ is $r+2$. Reverse its coefficients by setting
$H^*(x)=x^{r+2}H(1/x)$. Directly from~\eqref{eq:H},
\begin{equation}\label{eq:reverse}
 H^*(x)=(1+x)^{r+2}
 +(x+2x^2)\bigl((1+x)^p+(1+x)^q\bigr)+x^2+2x^3.
\end{equation}
By symmetry assume $p\leq q$ and set $u=q-p\geq0$. Define
\begin{equation}\label{eq:nested}
\begin{aligned}
 K_p(x)&=(1+x)^{p+2}+x+2x^2,\\
 R_{p,u}(x)&=(1+x)^uK_p(x)+x+2x^2,\\
 H^*(x)&=(1+x)^pR_{p,u}(x)+x^2+2x^3.
\end{aligned}
\end{equation}
The last identity follows by expanding the preceding two and comparing
with~\eqref{eq:reverse}.

\begin{lemma}[Convolution closure]\label{lem:convolution}
The product of two polynomials with nonnegative log-concave coefficient
sequences and no internal zeros again has these properties.
\end{lemma}

\begin{proof}
We recall the order-two Cauchy--Binet argument for convolution closure
\cite[Section~2]{Gross2015}. Extend each sequence by zeros outside its
support. For a sequence $c$, form the upper triangular Toeplitz matrix
$M(c)_{ij}=c_{j-i}$, $i,j\geq0$.
The ratios $c_{t+1}/c_t$ are nonincreasing on its positive support. It follows
that, for positive integers $A,B$ and any integer $m$,
\[
 c_{m+A}c_{m+B}\geq c_m c_{m+A+B}.
\]
Indeed, when the right side is positive all intervening terms are positive,
and comparison of $A$ successive ratios proves the inequality. When it is
zero the inequality is immediate. These are exactly the order-two minors
of $M(c)$; conversely the adjacent minors give log-concavity.
For convolution $c=a*b$ one has $M(c)=M(a)M(b)$.
Every relevant entry and minor can be computed in a finite truncation, so
Cauchy--Binet writes each order-two minor of $M(c)$ as a sum of products of
nonnegative minors. Finally, convolution of two positive interval supports
again has positive interval support.
\end{proof}

For fixed $h\geq0$ let
\[
 C(z,h)=\frac{z(z-1)\cdots(z-h+1)}{h!},\qquad C(z,0)=1,
\]
and set $C(z,h)=0$ for $h<0$. For a nonnegative integer $z$, these are
the binomial coefficients, including zero when $h>z$.
Writing $k_h,r_h,s_h$ for the coefficients of $K_p,R_{p,u},H^*$,
respectively,~\eqref{eq:nested} gives
\begin{equation}\label{eq:coefficients}
\begin{aligned}
 k_h&=C(p+2,h)+\ind_{h=1}+2\ind_{h=2},\\
 r_h&=\sum_{i=0}^h C(u,i)k_{h-i}+\ind_{h=1}+2\ind_{h=2},\\
 s_h&=\sum_{i=0}^h C(p,i)r_{h-i}+\ind_{h=2}+2\ind_{h=3}.
\end{aligned}
\end{equation}
Here each indicator is one when its subscripted condition holds and zero
otherwise. Define the differences, with zero coefficients outside support, by
\[
 \Delta^K_h=k_h^2-k_{h-1}k_{h+1},\quad
 \Delta^R_h=r_h^2-r_{h-1}r_{h+1},\quad
 \Delta^S_h=s_h^2-s_{h-1}s_{h+1}.
\]

\begin{lemma}\label{lem:certificates}
After substituting $p=a+1$, each of
\[
 \Delta^K_1,\Delta^K_2,\Delta^K_3,\quad
 \Delta^R_1,\Delta^R_2,\Delta^R_3,\quad
 \Delta^S_1,\Delta^S_2,\Delta^S_3,\Delta^S_4
\]
is a polynomial in $a,u$ with nonnegative rational coefficients.
\end{lemma}

\begin{proof}
Substitution and multiplication in~\eqref{eq:coefficients} give the
ten explicit expansions in Appendix~\ref{app:certificates}. All displayed
coefficients are nonnegative and the denominators are positive.
The identities involve only falling factorials through degree five, so they
are finite polynomial identities over $\mathbb{Q}[a,u]$.
\end{proof}

\begin{proof}[Proof of Theorem~\ref{thm:main}]
Every binomial row is log-concave and has positive interval support.
Adding $x+2x^2$ can alter a log-concavity difference only at centers $1,2,3$;
adding $x^2+2x^3$ can alter one only at centers $1,2,3,4$.
For instance, a difference uses only the coefficients at its center and
its two neighbors. Center zero is unchanged because the coefficient at
index $-1$ is zero.

Start with $(1+x)^{p+2}$. The three $K$ certificates prove that $K_p$ is
log-concave. Lemma~\ref{lem:convolution} then makes $(1+x)^uK_p$
log-concave, and the three $R$ certificates establish the same property
for $R_{p,u}$. Apply the convolution lemma once more, then the four $S$
certificates, to conclude that $H^*$ is log-concave.
Each intermediate polynomial has positive coefficients throughout its
support: its binomial part is positive there, and every sparse addition
lies inside that support. The argument includes $u=0$ and $p=1$.

Reversal preserves all log-concavity inequalities, so $H$ is log-concave.
Multiplying by $x^2$ gives $d_0=d_1=0$ and $d_2=H(0)>0$.
The inequalities at $j=1,2$ are therefore immediate; the remaining centers
$3\leq j\leq p+q+3$ are exactly the internal differences of $H$.
\end{proof}

\section{Exact verification and scope}

The symbolic identities above provide the quantified proof. Their transcription
was checked using exact rational arithmetic in SymPy~\cite{Meurer2017}.
The script \texttt{generate\_certificates.py} reconstructs
all ten differences directly from~\eqref{eq:coefficients}; its
\texttt{--check} option verifies the accompanying integer arrays and the
generated table source. All certificate entries are printed below, so the
mathematical argument does not require trusting a saved program verdict.

Additional integer checks compared~\eqref{eq:H} with a three-state tree
dynamic program for all $400$ ordered pairs $1\leq p,q\leq20$ and with
direct vertex-subset enumeration for all $36$ ordered pairs with
$p+q+4\leq13$. The states distinguish a selected root, an unselected root
dominated by a child, and an unselected root requiring its parent; their
recurrences are given in Du--Heilman--Panova~\cite[Section~4.1]{Du2026}.
A further check evaluated all $1{,}744{,}050$ differences for the
$11{,}325$ pairs $1\leq p\leq q\leq150$. No negative difference occurred.
These finite comparisons check the formula and implementation; they do not
establish an unbounded parameter claim.

The theorem concerns one block of exactly two bare vertices between two
leaf-bearing endpoints. Its nested decomposition supplies neither an
induction on arbitrary block length nor a treatment of interactions between
separated zero blocks. Thus it does not resolve the original universal
caterpillar question.

\paragraph{Computational and AI assistance.}
OpenAI language models assisted with proof development, exact verification
code, literature screening, and manuscript preparation. A separate model
context reconstructed the mathematical argument and performed additional
exact checks; this was an internal verification pass, not external peer
review. The reported computations used Python integers and rational symbolic arithmetic,
with no numerical root tests or floating-point sign decisions. The proof is
the displayed endpoint decomposition, convolution argument, and explicit
polynomial certificates.

\appendix
\section{The ten coefficient certificates}\label{app:certificates}

Throughout this appendix $p=a+1$ and $u=q-p$, so $a,u\geq0$ on the
parameter range of the theorem. Five of the expansions are short enough
to display directly:
\begin{align*}
 2\Delta^K_1&=a^2+11a+22,\\
 12\Delta^K_2&=a^4+10a^3+65a^2+200a+252,\\
 144\Delta^K_3&=(a+1)(a+2)(a+3)(a^3+9a^2+14a+24),\\
 2\Delta^R_1&=a^2+2au+15a+u^2+13u+36,\\
 2\Delta^S_1&=4a^2+4au+30a+u^2+15u+46.
\end{align*}
For each remaining difference, the array $C^{A,h}$ below means
\[
  b_h\Delta^A_h(a,u)=\sum_{i,j\geq0}C^{A,h}_{ij}a^iu^j,
  \qquad b_2=12,\quad b_3=144,\quad b_4=2880.
\]
Rows are indexed by the exponent $i$ of $a$, columns by the exponent $j$
of $u$, and every blank or unlisted entry is zero. These arrays and the
five formulas above specify all ten certificates completely. They can be
reconstructed by ordinary polynomial multiplication from
\eqref{eq:coefficients}; no parameter sampling is involved.

% Generated by generate_certificates.py; all entries are exact integers.
\begin{center}
\begin{minipage}{1\textwidth}
\centering\small
$C^{R,2}$: coefficients of $12\Delta^{R}_{2}$ in $a^i u^j$.\par\smallskip
\begin{tabular}{r|rrrrr}
$i\backslash j$ & 0 & 1 & 2 & 3 & 4 \\
\hline
0 & 528 & 376 & 101 & 14 & 1 \\
1 & 298 & 196 & 42 & 4 &  \\
2 & 77 & 36 & 6 &  &  \\
3 & 8 & 4 &  &  &  \\
4 & 1 &  &  &  &  \\
\end{tabular}
\end{minipage}
\end{center}
\begin{center}
\begin{minipage}{1\textwidth}
\centering\small
$C^{R,3}$: coefficients of $144\Delta^{R}_{3}$ in $a^i u^j$.\par\smallskip
\begin{tabular}{r|rrrrrrr}
$i\backslash j$ & 0 & 1 & 2 & 3 & 4 & 5 & 6 \\
\hline
0 & 144 & 1380 & 1276 & 615 & 163 & 21 & 1 \\
1 & 276 & 1436 & 1335 & 532 & 99 & 6 &  \\
2 & 220 & 837 & 618 & 174 & 15 &  &  \\
3 & 141 & 304 & 150 & 20 &  &  &  \\
4 & 67 & 69 & 15 &  &  &  &  \\
5 & 15 & 6 &  &  &  &  &  \\
6 & 1 &  &  &  &  &  &  \\
\end{tabular}
\end{minipage}
\end{center}
\begin{center}
\begin{minipage}{1\textwidth}
\centering\small
$C^{S,2}$: coefficients of $12\Delta^{S}_{2}$ in $a^i u^j$.\par\smallskip
\begin{tabular}{r|rrrrr}
$i\backslash j$ & 0 & 1 & 2 & 3 & 4 \\
\hline
0 & 1308 & 792 & 173 & 18 & 1 \\
1 & 1584 & 764 & 120 & 8 &  \\
2 & 764 & 252 & 24 &  &  \\
3 & 168 & 32 &  &  &  \\
4 & 16 &  &  &  &  \\
\end{tabular}
\end{minipage}
\end{center}
\begin{center}
\begin{minipage}{1\textwidth}
\centering\small
$C^{S,3}$: coefficients of $144\Delta^{S}_{3}$ in $a^i u^j$.\par\smallskip
\begin{tabular}{r|rrrrrrr}
$i\backslash j$ & 0 & 1 & 2 & 3 & 4 & 5 & 6 \\
\hline
0 & 12528 & 13572 & 5920 & 1521 & 271 & 27 & 1 \\
1 & 27144 & 24544 & 9390 & 2060 & 258 & 12 &  \\
2 & 24544 & 19044 & 6108 & 996 & 60 &  &  \\
3 & 12696 & 8096 & 1944 & 160 &  &  &  \\
4 & 4048 & 1920 & 240 &  &  &  &  \\
5 & 768 & 192 &  &  &  &  &  \\
6 & 64 &  &  &  &  &  &  \\
\end{tabular}
\end{minipage}
\end{center}
\begin{center}
\begin{minipage}{1\textwidth}
\centering\small
$C^{S,4}$: coefficients of $2880\Delta^{S}_{4}$ in $a^i u^j$.\par\smallskip
\begin{tabular}{r|rrrrrrrrr}
$i\backslash j$ & 0 & 1 & 2 & 3 & 4 & 5 & 6 & 7 & 8 \\
\hline
0 & 2880 & 18240 & 28844 & 21968 & 11049 & 2960 & 426 & 32 & 1 \\
1 & 36480 & 101456 & 115144 & 71832 & 24920 & 4632 & 432 & 16 &  \\
2 & 101456 & 213624 & 191736 & 87600 & 21000 & 2472 & 112 &  &  \\
3 & 142416 & 239808 & 160400 & 51600 & 7840 & 448 &  &  &  \\
4 & 119904 & 153000 & 72960 & 15000 & 1120 &  &  &  &  \\
5 & 61200 & 56592 & 17424 & 1792 &  &  &  &  &  \\
6 & 18864 & 11424 & 1792 &  &  &  &  &  &  \\
7 & 3264 & 1024 &  &  &  &  &  &  &  \\
8 & 256 &  &  &  &  &  &  &  &  \\
\end{tabular}
\end{minipage}
\end{center}

\clearpage
\bibliographystyle{plainnat}
\bibliography{references}
\end{document}
